\section{Benchmark Metrics}
\label{sec:metrics}

Let $Q$ be the $n$ frozen questions of a submission, with outcome labels
$o(q)$ and premise labels $p(q)$ as in Table~\ref{tab:taxonomy}, and let
$E(q)$ be the set of independent post-cutoff papers cited as supporting
evidence for $q$'s label. All rates are over $n$, so the four outcome
rates sum to one.

\begin{table}[t]
  \centering
  \small
  \caption{Benchmark metrics v1.0. All are computed by
  \texttt{benchmark/metrics.py} from the released annotation records.}
  \label{tab:metrics}
  \begin{tabular}{@{}lp{9.6cm}@{}}
    \toprule
    Metric & Definition \\
    \midrule
    Coverage (future attention rate) &
      $\frac{1}{n}\,|\{q : o(q) \neq \labelfmt{not\_addressed}\}|$ ---
      did future science engage the question at all? \\
    Answer rate &
      $\frac{1}{n}\,|\{q : o(q) = \labelfmt{answered}\}|$ \\
    Partial rate &
      $\frac{1}{n}\,|\{q : o(q) = \labelfmt{partially\_addressed}\}|$ \\
    Open rate &
      $\frac{1}{n}\,|\{q : o(q) = \labelfmt{posed\_but\_open}\}|$ ---
      the community independently recognized the question \\
    Premise refutation rate &
      $\frac{1}{n}\,|\{q : p(q) = \labelfmt{refuted}\}|$ ---
      questions that led to an established conclusion being overturned \\
    Evidence strength &
      mean $|E(q)|$ over engaged questions, and the multi-source rate
      $\frac{1}{n}\,|\{q : |E(q)| \geq 2\}|$ --- is the label supported
      by multiple independent papers? \\
    Lead time &
      years between question submission and the first time the community
      \emph{independently poses} the same question \\
    Community attention &
      volume of future investment engaging the question: papers,
      citations, review mentions, and major observing programs (e.g.\
      JWST/HST proposals) \\
    \bottomrule
  \end{tabular}
\end{table}

\paragraph{Coverage vs.\ answer rate.} Coverage asks whether the
question pointed anywhere the community went at all; the answer/partial/
open decomposition asks what happened when it got there. A system can
maximize coverage with fashionable-topic questions, which is why
coverage is never reported alone (Section~\ref{sec:threats} discusses
the popularity confound).

\paragraph{Premise refutation rate.} This is the metric we most want the
field to adopt. Questions that trigger refutations are the rarest and
arguably most valuable output of question discovery---they mark places
where the literature's accepted conclusions were wrong and where a
well-aimed question preceded the correction. Under the two-dimensional
taxonomy the refutation is recorded without erasing the fact that the
question was thereby \emph{answered}.

\paragraph{Lead time.} If a system poses a question at the cutoff and
the community first independently poses it in year $T{+}\ell$, the
system led the field by $\ell$ years; averaged over questions this
yields a comparable earliness score. Measuring $\ell$ requires
identifying \emph{community first-posed dates}, which demands careful
review-literature annotation we do not yet have. Astronomy~v1 therefore
reports \texttt{mean\_lead\_time\_years:\,null} rather than a number we
cannot defend; the released records do include a weaker, well-defined
lower bound---\emph{first-engagement lag}, the years from cutoff to the
earliest judge-cited supporting paper (pilot mean 2.9, range 1--5)---
which should not be confused with lead time.

\paragraph{Community attention.} Beyond binary engagement, the volume of
future investment (paper counts, citations to engaging papers, review
mentions, dedicated observing programs) reflects how much the community
cared. v1 records the ingredients (supporting bibcodes, their venues and
years) and reports evidence strength; a calibrated attention index is
future work, and Section~\ref{sec:threats} explains why raw attention
must never be the headline metric.

\paragraph{Reporting requirements.} A benchmark report must state: the
instance and protocol versions, $n$, all outcome and premise rates, the
evidence-strength pair, and either lead time or an explicit null. The
released implementation produces exactly this
(\texttt{results/astronomy\_v1/metrics.json}) with one command, and CI
fails if the committed numbers do not reproduce from the raw
annotations.
